\section{Design Validation}\label{sec:test}

\subsection{Testing Materials}
Most of the testing planned will be performed with the device and simple materials that can be found in the engineering facilities,
for example scales or various weights.

Part of the testing of this device is intended to be performed attached to a drone. There is a drone available in the engineering facilities and ideally it will be used for testing the device.
One goal of this device is to be simple to use without special adaptations to a drones hardware or control software aside from attaching, powering, and commanding the device,
so using an off the shelf drone is a good test case. However there are a few possible technical difficulties in using such a drone for this test:
\begin{itemize}
    \item The drone available is not made to carry additional weight, and may not be capable of lifting the device.
    \item The device will require 12V power, which will be difficult to get from the off the shelf drone, 
        and adding a secondary battery will add significant weight.
\end{itemize}

If using the off the shelf drone proves infeasible, then a simple drone may need to be built for the test. The author has most of the necessary components listed in table \ref{tab:drone},
and a simple frame can be fabricated in the engineering facilities. This device need not be durable or particularly reliable for these tests.

\begin{table}[ht]
    \centering
    \caption{Estimated budget for a test drone.}
    \label{tab:drone}

    \begin{spreadtab}{{tabular}{|c|c|c|c||c|}}
        \hline
        @Item                       & @Supplier                                                                                                         & @Unit Cost (USD) & @Quantity & @Total Cost (USD) \\
        \hline
        @Motors                     & @-                                                                                                                &     0 & 4 & [-1,0] * [-2,0] tag(start) \\
        @Motor Controllers          & @-                                                                                                                &     0 & 4 & [-1,0] * [-2,0] \\
        @3 Cell LiPo                & @-                                                                                                                &     0 & 1 & [-1,0] * [-2,0] \\
        @PX4 Flight Controller      & @-                                                                                                                &     0 & 1 & [-1,0] * [-2,0] \\
        @RC Radio and Remote        & @-                                                                                                                &     0 & 1 & [-1,0] * [-2,0] \\
        @Telemetry Radio            & @-                                                                                                                &     0 & 1 & [-1,0] * [-2,0] \\
        @Clockwise Propellers       & @\href{https://www.robotshop.com/products/8-x-45-orange-multi-rotor-cw-propeller-pair}{RobotShop}                 &  3.78 & 3 & [-1,0] * [-2,0] \\  
        @Counterclockwise Propellers& @\href{https://www.robotshop.com/products/8-x-45-orange-multi-rotor-ccw-propeller-pair}{RobotShop}                &  3.78 & 3 & [-1,0] * [-2,0] \\    
        @RobotShop Shipping         & @-                                                                                                                &  4.74 & 1 & [-1,0] * [-2,0] tag(pretax)\\
        @Tax                        & @- & sum(cell(start):cell(pretax)) * 0.07 & 1 & [-1,0] * [-2,0] tag(end) \\
        \hline
        \multicolumn{4}{|c||}{@Total} & sum(cell(start):cell(end)) \\
        \hline
    \end{spreadtab}
\end{table}

\subsection{Attaching and Detaching}

If a drone is available, it will be used here. Otherwise the device will be positioned beneath a perch using string, a pedestal,
or similar means. The following procedure will be performed:
\begin{enumerate}
    \item The device will begin in its "home" position.
    \item The device will be electronically triggered to extend and grasp the perch.
    \item Within a period of 2 seconds, the device should have grasped the perch and engaged any locking mechanisms.
    \item The supports used to position the device will be removed. If a drone is used, its throttle will be set to 0.
    \item The device should support the weight of itself, and a drone if present, without disengaging.
\end{enumerate}

Once this is complete, the following procedure will be performed to test detaching:
\begin{enumerate}
    \item If a drone is being used, its throttle will be brought to slightly below the level that was previously used to attach.
    \item The device will be electronically triggered to detach.
    \item Within 2 seconds, the device should have released any locking mechanisms and freed from its perch.
    \item The grasper should return to its home positions.
    \item If a drone is used, it should be in stable flight free from the perch. This may be unattainable due to the skill of the pilot,
            rather than the design.
\end{enumerate}

These two procedures will be repeated for positions 5cm away from ideal in all directions.


\subsection{Remaining Attached}

To verify the holding strength of the device, it will first be attached to a perch with a 500g mass attached to the device
to simulate the drones mass.
To simulate the dynamic effects of wind, an additional weight of 3.22 kg will be added. This supplies a force equivalent
to the maximum dynamic forces estimated due to swaying of the perch.

The device will pass this test if it remains perched for 10 minutes %?
without detaching or visible damage to the device.

If the design is determined to be unable of detaching unintentionally without critical failure of the device 
itself (for example, the "lock" like concept), then testing will not be repeated for different surfaces or perch shapes.

\subsection{Perch Identification}

To test the perch identification method, the device will be placed a known distance away from a suitable
perch. The measured distance to the perch will be collected from the device, and the error computed
against the known distance. 

The error must be within the limits of the grasping mechanisms tolerance; for example if the grasper
is able to grasp a perch that is up to 10mm misaligned from its center then the allowed error of the
detection method is 10mm.


\subsection{Size}

The device will be measured in its home position, and the measurements will be compared to those of a commercial drone 
to confirm that, with the device mounted at a reasonable position, the drones bounding box is not extended past the specified
limits.

\subsection{Weight}

The weight of the final device prototype will be verified with the tools available in the school of engineering facilities.


\pagebreak
